\section{Pre-registered Tests: What Held and What Did Not}
\label{sec:outcomes}

Table~\ref{tab:prereg} lists the confirmatory tests of the amendments that bear on the claims above (full log: Supp.~A).
Predictions that failed or could not be evaluated include the image benefit on text-undecided questions in our fine-tuned models (A9), a competition score (A15), a median-margin rule for the image gap (A16), the text margin as a predictor of the dose-response (A17), three attempts to remove the cue from Med-CMR (A6, A7, A12), cue uptake and strong lock-out on A-OKVQA (A21), the boundary shift of the position shortcut (A22) and three methods on the real shortcut (A24); a saturation-aware loss and an image-evidence error alarm did not help either (A12).

\begin{table*}[t]
  \centering
  \footnotesize
  \setlength{\tabcolsep}{4pt}
  \begin{tabular}{@{}l p{0.37\textwidth} p{0.50\textwidth}@{}}
    \toprule
    Amend. & Prediction & Outcome \\
    \midrule
    A8 & Text-decisive answers never change; the image is a tie-breaker on undecided questions (released adapter) & all six predictions supported (0 of 2{,}463 changed; $+12.1$ points, $p=2\times10^{-5}$) \\
    A9--A10 & The tie-breaker benefit ($\ge+5$ points) replicates in fine-tuned Qwen 3B/7B and in four recipe variants & \emph{falsified}: $-1.0$, $+1.8$ and $-4.1$ to $+2.6$ points; no-veto, habit and alarm supported \\
    A11 & PMC-VQA (no cue): no habit, gap $\ge+5$, undecided-stratum effect $\ge+5$ (Q1--Q3), alarm (Q5) & Q1--Q3 supported for three families; Q4 not evaluable (empty band); Q5 supported on E1, not on E2 \\
    A13 & Cue planted at 92\%: habit, gap $\approx0$, text margin $\ge0.99$ for 30\% of questions (I1--I4) & I1, I3, I4 supported; I2 \emph{falsified} (confident but never $\ge0.99$) \\
    A15 & Competition score beats saturation (C1, C2); blur factorial (F1--F3) & C1, C2 \emph{falsified}; F1--F3 supported for three families \\
    A16 & 7B seeds (S1, S2); OmniMedVQA (O1, O3, O4) and the lock rule (O2) & S1, S2, O1, O3, O4 supported; O2 \emph{falsified} for three systems \\
    A17 & Dose-response H1, H2, H3a, H4, H5; position shortcut P1--P4; H3b & H1, H2, H3a, H4, H5 and P1--P4 supported (Holm); H3b \emph{falsified} \\
    A18 & Rescue: M-A, M-B, M-D, M-E, M-G (controlled setting); M-F (Med-CMR) & M-A, M-B, M-D, M-E, M-G supported; M-F \emph{falsified} (collapse to chance) \\
    A19 & Held-out transfer to InternVL3-2B and Qwen-7B (A, B, C) & \emph{supported} for both models: deviations $0.65$ and $1.96$ points ($\le4$), boundary and monotonicity as predicted \\
    A21 & Non-medical replication on A-OKVQA: dose-response (R1), cue uptake (R2), boundary shift (R3), lock-out at $d=0.92$ (R4) & R1, R3 supported in three series; R2 \emph{falsified} (accuracy on cue-free questions does not fall); R4 \emph{falsified} ($+8$ to $+12$ points at $d=0.92$) \\
    A22 & Higher strengths (H1--H3) and position shortcut (P1--P4) on A-OKVQA; second seed of A19 (S) & H1, H2 supported, lock-out at $d=0.99$; P1, P2 supported, P3 \emph{falsified}; S supported for both models \\
    A23 & Second seed of the A-OKVQA lock-out (B1, B2); Qwen-7B on A-OKVQA (C1--C3) & all supported: lock-out at $d=0.99$ in a second seed ($+0.09$ points) and in Qwen-7B ($+0.09$) \\
    A24 & Real Med-CMR shortcut: explicit length prior (LP), mismatched-image uniformity (MU), both; criterion E1 & E1 \emph{falsified} for all three (image effect $\le+4.6$ points; accuracy $79$--$80\%$ for LP+MU) \\
    A12, A6--A7 & Cue-removal tests; alarm combiner (M1); saturation-aware loss (M2) & L1 not evaluable; A6, A7 not evaluable; M1, M2 \emph{falsified} \\
    \bottomrule
  \end{tabular}
  \caption{\textbf{Confirmatory pre-registered tests and outcomes.} ``Supported'' means every stated criterion held; thresholds were fixed before the runs (hash-chained amendments).}
  \label{tab:prereg}
\end{table*}
